% TEMPLATE-MILC-v3.tex - plantilla maestra UNICA de las guias MILC v3.
%
% La usan las cuatro familias de guias, cada una con su driver:
%   · Grados 6-11          -> scripts/build-guias-g11.py
%   · Bebras               -> scripts/build-guias-bebras.py
%   · Semillero            -> scripts/build-guias-semillero.py
%   · Territorio interior  -> scripts/build-guias-territorio-interior.py
%
% Antes habia tres copias casi identicas (~400 lineas iguales, 6 distintas):
% cada mejora habia que aplicarla tres veces, y en la practica no se hacia
% (el motor de recursos vivia solo en dos de ellas). Lo que varia entre
% familias esta parametrizado en 6 placeholders que cada driver llena
% (se nombran aqui SIN su sintaxis real: el builder reemplaza por texto plano,
% tambien dentro de los comentarios, y un valor multilinea rompe el preambulo):
%
%   HEADER_IZQ        encabezado izquierdo de pagina
%   HEADER_DER        encabezado derecho de pagina
%   PORTADA_ETIQUETA  rotulo sobre el numero grande (GRADO/MOMENTO/MODULO)
%   PORTADA_SERIE     linea de serie de la portada
%   PORTADA_ID        identificador de la guia en la portada
%   PORTADA_CIRCULO   el circulito de grado (°), o vacio si no aplica
%   PDF_TITULO        titulo en texto plano (sin \textbf ni comillas TeX) para
%                     los metadatos del PDF (/Title); lo produce lib_xelatex.texto_pdf
%
% Composicion (auditoria 2026-09-03): las bandas de estacion piden espacio
% (\Needspace*) y prohiben el salto justo despues (\nopagebreak), asi no quedan
% huerfanas al pie; las cajas largas (softbox, drawbox, infoband, puentebox) son
% `breakable`, asi una caja de media pagina ya no arrastra todo a la siguiente
% dejando la anterior al 40 %. Todo texto sobre mostaza o verde va en milcNegro
% (blanco sobre #E5B400 da 1,93:1 y sobre #58A923 2,95:1; ninguno pasa WCAG AA).
% El resto de claves las documenta content/guias/_SCHEMA.md. El script valida
% que no queden placeholders sin reemplazar antes de escribir el archivo final.

% Etiquetado PDF (latex-lab): /Lang, estructura y texto alternativo de las
% figuras, para lectores de pantalla. Debe ir ANTES de \documentclass.
\DocumentMetadata{lang=es-CO,tagging=on}
\documentclass[11pt,letterpaper]{article}
% headheight=24pt: la cabecera derecha lleva dos lineas (identidad de la guia
% y estacion actual). headsep se acorta en la misma medida para que el bloque
% de texto quede exactamente donde estaba con headheight=12pt.
\usepackage[margin=2.0cm,headheight=24pt,headsep=13pt]{geometry}
\usepackage[table]{xcolor}
\usepackage{fontspec}
\usepackage[spanish]{babel}
\usepackage{tikz}
% Librerías TikZ para los diagramas de las guías (bloque `recursos:`):
%  · babel  → babel en español vuelve activos caracteres como «>», y TikZ los usa
%             en sus opciones (>=stealth, ->). Sin esta librería, cualquier
%             diagrama con flechas revienta con «\language@active@arg> has an
%             extra }». Debe ir ANTES de que se dibuje nada.
%  · positioning        → sintaxis «below left=2pt and 3pt».
%  · decorations.pathreplacing → llaves ({brace}) para señalar rangos.
\usetikzlibrary{babel,positioning,decorations.pathreplacing}
\usepackage{graphicx}
% Circuitos: el trabajo de electrónica se hace hoy con micro:bit y sus sensores
% integrados (MakeCode, sin cableado), así que no se necesitan esquemáticos.
% Si algún día entran montajes con protoboard, basta descomentar esta línea:
% los diagramas .tex/.tikz declarados en `recursos:` ya se incrustan solos.
% \usepackage{circuitikz}
\usepackage[most]{tcolorbox}
\usepackage{tabularx}
\usepackage{array}
\usepackage{fancyhdr}
\usepackage{titlesec}
\usepackage[document]{ragged2e}  % bandera a la derecha en todo el cuerpo
\usepackage{enumitem}
\usepackage{microtype}
\usepackage{etoolbox}
\usepackage{needspace}
% hyperref al final del preambulo: metadatos del PDF (titulo, autor, idioma,
% familia), marcadores por estacion (\pdfbookmark) y enlaces sin recuadro.
\usepackage[hidelinks,
  pdftitle={Audio — grabar, limpiar y exportar},
  pdfauthor={PhD. Álvaro Cárdenas Orozco},
  pdfsubject={Tecnología e Informática · Grado 8 · Guía 25 · MILC v3},
  pdflang={es-CO},
  bookmarksopen=true,bookmarksnumbered=false]{hyperref}

% Helvetica Neue no trae la flecha U+2192, asi que cada «→» escrito literalmente
% en el YAML se compilaba a NADA, en silencio (462 flechas en 61 guias: rutas del
% tipo «paleta → tipografia → lockup» salian rotas). Mapeamos el caracter a su
% version matematica, que si tiene glifo. Asi el autor puede seguir escribiendo
% «→» comodamente en el YAML.
\usepackage{newunicodechar}
\newunicodechar{→}{\ensuremath{\rightarrow}}
% Mismo problema con los simbolos de marca y vinetas: el «✓» de «una palomita ✓»
% salia como cuadrito vacio en el PDF de todas las guias que lo usaban.
\usepackage{pifont}
\newunicodechar{✓}{\ding{51}}
\newunicodechar{✗}{\ding{55}}
\newunicodechar{✔}{\ding{52}}
\newunicodechar{•}{\textbullet}
\newunicodechar{≥}{\ensuremath{\geq}}
\newunicodechar{≤}{\ensuremath{\leq}}

\setmainfont{Helvetica Neue}
\setsansfont{Helvetica Neue}
\setlength{\parindent}{0pt}
\setlength{\parskip}{6pt}
% Sin particion de palabras: en 6.º-7.º una palabra cortada por guion al final
% de linea («Comuni-cacion») frena la lectura mucho mas que una linea corta.
\hyphenpenalty=10000
\exhyphenpenalty=10000
\pretolerance=2000
\tolerance=3000
\setlength{\emergencystretch}{3em}
\linespread{1.10}
\renewcommand{\arraystretch}{1.25}
\setlist{leftmargin=5mm,itemsep=1mm,topsep=1mm}
\newcolumntype{Y}{>{\RaggedRight\arraybackslash}X}

\definecolor{milcMagenta}{HTML}{E6007E}
\definecolor{milcVino}{HTML}{5A0038}
\definecolor{milcTurquesa}{HTML}{0093A5}
\definecolor{milcVerde}{HTML}{58A923}
\definecolor{milcMostaza}{HTML}{E5B400}
\definecolor{milcOcre}{HTML}{B86600}
\definecolor{milcNegro}{HTML}{241816}
\definecolor{milcGris}{HTML}{F7F7F4}
\definecolor{milcLinea}{HTML}{E8E2DD}
\definecolor{milcGradePrimary}{HTML}{FF6600}
\definecolor{milcGradeDark}{HTML}{9A3E00}
\definecolor{milcGradeSoft}{HTML}{FFE4D0}
\definecolor{milcGradeAccent}{HTML}{CFE7FF}
\definecolor{milcGradeText}{HTML}{FFFFFF}
\color{milcNegro}

\pagestyle{fancy}
\fancyhf{}
% Cabecera de dos lineas: identidad de la guia arriba y, a la derecha y en
% gris, la estacion en curso (\markright desde \milcestacion). Antes de la
% primera estacion la segunda linea va vacia; el \strut mantiene la altura
% para que la linea de arriba no baile entre paginas.
\lhead{\small Tecnología e Informática · Grado 8\\[-1pt]{\footnotesize\strut}}
\rhead{\small Guía 25 · MILC v3\\[-1pt]{\footnotesize\strut\color{milcNegro!65}\rightmark}}
\cfoot{\small\thepage}
\renewcommand{\headrulewidth}{0pt}

% Sobre mostaza (#E5B400) y verde (#58A923) el texto va en milcNegro; sobre el
% resto de la paleta, en blanco. Se usa dentro de fonttitle/fontupper, donde un
% \color posterior manda sobre coltitle.
\newcommand{\milctintasobre}[1]{%
  \ifboolexpr{test{\ifstrequal{#1}{milcMostaza}} or test{\ifstrequal{#1}{milcVerde}}}%
    {\color{milcNegro}}{\color{white}}}

\newtcolorbox{titlebox}[1]{
  enhanced,colback=#1,colframe=#1,arc=7mm,boxrule=0pt,
  left=7mm,right=7mm,top=3mm,bottom=3mm,
  before skip=8pt,after skip=8pt,
  fontupper=\sffamily\bfseries\Large\color{white}
}

\newtcolorbox{softbox}[3]{
  enhanced,breakable,pad at break=3mm,
  colback=#2,colframe=#1,borderline west={4pt}{0pt}{#1},
  arc=2mm,boxrule=.4pt,left=4mm,right=4mm,top=3mm,bottom=3mm,
  before skip=6pt,after skip=8pt,title={#3},coltitle=white,
  colbacktitle=#1,fonttitle=\sffamily\bfseries\milctintasobre{#1},fontupper=\RaggedRight,
  attach boxed title to top left={xshift=3mm,yshift=-2mm},
  boxed title style={arc=2mm, boxrule=0pt}
}

\newtcolorbox{drawbox}[2]{
  enhanced,breakable,pad at break=4mm,
  colback=white,colframe=#1,arc=4mm,boxrule=1pt,
  left=5mm,right=5mm,top=4mm,bottom=4mm,before skip=8pt,after skip=8pt,
  title={#2},coltitle=white,colbacktitle=#1,fonttitle=\sffamily\bfseries\milctintasobre{#1},
  fontupper=\RaggedRight,
  attach boxed title to top left={xshift=4mm,yshift=-2mm},
  boxed title style={arc=3mm, boxrule=0pt}
}

\newtcolorbox{infoband}[1]{
  enhanced,breakable,pad at break=2mm,colback=#1!8,colframe=#1!50,arc=2mm,boxrule=.5pt,
  left=4mm,right=4mm,top=2mm,bottom=2mm,before skip=4pt,after skip=4pt,
  fontupper=\small\RaggedRight
}

% Puente narrativo entre secciones (contrato editorial Conéctate).
% Da continuidad: "Acabas de X. Ahora vas a Y porque Z."
% Visualmente: banda izquierda en color del grado, texto en itálica pequeña.
\newtcolorbox{puentebox}{
  enhanced,breakable,pad at break=2mm,colback=milcGradeSoft,colframe=milcGradeDark,
  borderline west={3pt}{0pt}{milcGradeDark},
  arc=0mm,boxrule=0pt,left=5mm,right=4mm,top=2.5mm,bottom=2.5mm,
  before skip=4pt,after skip=8pt,
  fontupper=\itshape\small\color{milcNegro}\RaggedRight
}
% La flecha va en modo matematico a proposito: Helvetica Neue NO tiene el glifo
% U+2192, asi que \textrightarrow se compilaba a NADA («Missing character: There
% is no → in font Helvetica Neue») y el puente salia sin flecha. En modo
% matematico la flecha viene de la fuente matematica y si se dibuja.
\newcommand{\puente}[1]{%
  \begin{puentebox}%
    \textcolor{milcGradeDark}{$\rightarrow$}\hspace{2mm}#1%
  \end{puentebox}%
}

\newcommand{\linead}[1]{\noindent\color{milcLinea}\rule{#1}{0.5pt}\color{milcNegro}}
\newcommand{\respuesta}[1]{\par\vspace{2mm}\foreach \n in {1,...,#1}{\noindent\color{milcLinea}\rule{\linewidth}{0.5pt}\par\vspace{5mm}}\color{milcNegro}}
\newcommand{\checkbox}{\textcolor{milcGradeDark}{\fbox{\phantom{x}}}\hspace{5pt}}

% ─── Listas aireadas para las guías socioemocionales ────────────────────
% Componer las verificaciones como «\checkbox texto\\» deja el texto que salta
% de línea debajo del cuadrito y apelmaza el bloque. Estos entornos alinean la
% continuación y airean los ítems. Los usa Territorio Interior; cualquier
% familia puede hacerlo.
\newlist{checklista}{itemize}{1}
\setlist[checklista]{
  label=\textcolor{milcGradeDark}{\fbox{\phantom{x}}},
  leftmargin=9mm, labelsep=3.2mm, itemsep=2.4mm,
  topsep=2.5mm, parsep=0pt, partopsep=0pt, before=\RaggedRight
}
\newlist{pasoslista}{enumerate}{1}
\setlist[pasoslista]{
  label=\textcolor{milcGradeDark}{\sffamily\bfseries\arabic*.},
  leftmargin=8mm, labelsep=3mm, itemsep=2.8mm,
  topsep=1mm, parsep=0pt, partopsep=0pt, before=\RaggedRight
}

\newcommand{\phasecard}[4]{
\begin{tcolorbox}[enhanced,colback=#3,colframe=#2,arc=3mm,boxrule=.5pt,height=3.05cm,valign=center,halign=center,left=1.5mm,right=1.5mm,top=2mm,bottom=2mm]
{\sffamily\bfseries\fontsize{22}{24}\selectfont\textcolor{#2}{#1}}\par
{\sffamily\bfseries\small #4}
\end{tcolorbox}
}

% ── Piezas del rediseno 2026 (legibilidad 6.º-11.º) ───────────────────
% Banda de estacion: reemplaza el titlebox macizo. Titulo a la izquierda,
% dato practico (tiempo/modalidad) a la derecha, en una sola linea.
% Nunca huerfana: pide 9 lineas libres (o salta de pagina rellenando con
% \vfill) y prohibe el salto justo despues de la banda. Nueve y no seis:
% la caja partible que sigue necesita ~80pt (titulo adosado, relleno y dos
% lineas) para arrancar en la misma pagina; con menos, tcolorbox la manda
% entera a la siguiente y la banda se queda sola. El penalty 10000 va
% en `after`, ANTES del espacio posterior: un `after skip` (aunque sea 0pt)
% es un \vskip tras la caja, y ese glue es un punto de corte legal que un
% \nopagebreak escrito despues de \end{tcolorbox} ya no cierra. Cada banda
% deja marcador PDF y actualiza la cabecera derecha.
\newcounter{milcestacion}
\newcommand{\milcestacion}[2]{%
\Needspace*{9\baselineskip}%
\stepcounter{milcestacion}%
\pdfbookmark[1]{#2}{milcestacion\themilcestacion}%
\markright{#2}%
\begin{tcolorbox}[enhanced,colback=#1,colframe=#1,arc=2mm,boxrule=0pt,
  left=5mm,right=5mm,top=2.6mm,bottom=2.6mm,before skip=11pt,
  after={\par\nopagebreak\vskip7pt}]
{\sffamily\bfseries\fontsize{15}{18}\selectfont\milctintasobre{#1}#2}
\end{tcolorbox}}

% Tarjeta de paso numerada: sustituye las tablas «Pilar / Aplicacion», que
% eran formato de consulta docente, no de estudiante. El numero va en un
% circulo sobre el borde para que la secuencia se lea de un vistazo.
\newcommand{\milcpaso}[3]{%
\Needspace{4\baselineskip}%
\begin{tcolorbox}[enhanced,colback=white,colframe=milcLinea,arc=1.5mm,boxrule=.9pt,
  left=14mm,right=4mm,top=3mm,bottom=3mm,before skip=4pt,after skip=4pt,
  fontupper=\RaggedRight,
  overlay={\node[anchor=north west,circle,fill=#2,text=white,inner sep=0pt,
    minimum size=7.4mm,font=\sffamily\bfseries\small]
    at ([xshift=3.6mm,yshift=-3.2mm]frame.north west) {#1};}]
#3
\end{tcolorbox}}

% Casilla marcable de tamano util con lapiz (la anterior era \fbox{\phantom{x}},
% de alto variable segun el texto que la seguia).
\newcommand{\milcmarca}{\framebox[3.8mm]{\rule{0pt}{3.4mm}}}

% Fila marcable de lista suelta (checks de la estacion 1).
\newcommand{\milccheckrow}[1]{%
\par\vspace{1.8mm}\noindent
\makebox[6.4mm][l]{\raisebox{-.55mm}{\textcolor{milcNegro}{\milcmarca}}}%
\parbox[t]{\dimexpr\linewidth-6.4mm\relax}{\RaggedRight #1}\par}

% Celda marcable dentro de la rubrica (tabla de autoevaluacion). Misma
% sangria francesa, pero pensada para vivir dentro de una columna X.
\newcommand{\milcopcion}[1]{%
\makebox[5.2mm][l]{\raisebox{-.55mm}{\textcolor{milcNegro}{\milcmarca}}}%
\parbox[t]{\dimexpr\linewidth-5.2mm\relax}{\RaggedRight #1}}

% ── Recursos visuales de la guía ──────────────────────────────────────
% Declarados en el bloque `recursos:` del YAML y emitidos por el builder.
% \guiaFigura   → imagen rasterizada o PDF (foto, carta celeste, montaje).
% \guiaDiagrama → fuente TikZ/circuitikz (.tex) incrustada como vector:
%                 es la vía para circuitos y esquemáticos editables.
% [#1] = texto alternativo (alt) · #2 = ruta relativa al .tex · #3 = pie
% El alt llega al PDF etiquetado (/Alt) y lo lee el lector de pantalla.
\newcommand{\guiaFigura}[3][]{%
\begin{center}
\includegraphics[width=0.88\linewidth,height=0.38\textheight,keepaspectratio,alt={#1}]{#2}\\[1.5mm]
{\footnotesize\itshape\textcolor{milcNegro!75}{#3}}
\end{center}
}

\newcommand{\guiaDiagrama}[2]{%
\begin{center}
\input{#1}\\[1.5mm]
{\footnotesize\itshape\textcolor{milcNegro!75}{#2}}
\end{center}
}

\begin{document}
\thispagestyle{empty}

% ── Portada ───────────────────────────────────────────────────────────
% Antes: un rectangulo de color a pagina completa. Se veia bien en pantalla,
% pero en la fotocopiadora del colegio se comia el toner y salia gris sucio.
% Ahora la banda ocupa el tercio superior y el resto va en blanco.
\begin{tikzpicture}[remember picture,overlay]
  \path[fill=milcGradePrimary]
    (current page.north west) rectangle ([yshift=-10.6cm]current page.north east);

  \node[anchor=north west,text=milcGradeText] at ([xshift=2.0cm,yshift=-1.55cm]current page.north west)
    {\sffamily\bfseries\fontsize{12}{14}\selectfont GRADO};
  \node[anchor=north west,text=milcGradeText,opacity=.85] at ([xshift=2.0cm,yshift=-2.35cm]current page.north west)
    {\sffamily\bfseries\fontsize{8.5}{10}\selectfont SERIE GUÍAS · TECNOLOGÍA E INFORMÁTICA};
  \node[anchor=north east,text=milcGradeText,opacity=.85,align=right] at ([xshift=-2.0cm,yshift=-1.55cm]current page.north east)
    {\sffamily\bfseries\fontsize{9}{12}\selectfont I.E. Sor María Juliana\\Cartago, Valle del Cauca};

  \node[anchor=west,text=milcGradeText] (gradonum) at ([xshift=1.85cm,yshift=-7.1cm]current page.north west)
    {\sffamily\bfseries\fontsize{112}{112}\selectfont 8};
  % El circulito de grado (°) solo aplica a las guias de grado («11°»); en
  % Bebras y Semillero el numero grande es un momento/modulo, no un grado.
  % Se posiciona relativo al nodo `gradonum`, no en coordenadas absolutas.
  \draw[milcGradeText,line width=1.6pt]
    ([xshift=2.0mm,yshift=-4.5mm]gradonum.north east) circle (.15cm);

  \node[anchor=west,text=milcGradeText,text width=11.4cm,align=left] at ([xshift=1.5cm,yshift=0.5cm]gradonum.east)
    {
     {\sffamily\bfseries\fontsize{9}{11}\selectfont\MakeUppercase{Período 3 · Multimedia, narrativa y ciberseguridad}}\\[3mm]
     {\sffamily\bfseries\fontsize{21}{25}\selectfont\setlength{\baselineskip}{27pt}Grabar, limpiar\\[4pt]y exportar}\\[3.5mm]
     {\sffamily\bfseries\fontsize{15}{18}\selectfont Guía 25-8-TIC}
    };
\end{tikzpicture}

\vspace*{9.4cm}
\vfill

{\sffamily\bfseries\fontsize{8}{10}\selectfont\textcolor{milcGradeDark}{LO QUE TE LLEVAS HOY}}

\begin{tcolorbox}[enhanced,colback=white,colframe=milcNegro,arc=2mm,boxrule=1.6pt,
  left=5mm,right=5mm,top=4mm,bottom=4mm,before skip=3pt,after skip=12pt,
  fontupper=\RaggedRight\fontsize{11}{15.5}\selectfont]
Una pieza de audio de uno a dos minutos, una entrevista corta a un familiar, una narración o un pódcast breve, grabada con el celular o el computador y editada con Audacity o una app del celular con tres operaciones: recorte, ajuste de volumen y limpieza de silencio o ruido. Exportada en dos formatos, MP3 y WAV, con una ficha de cinco líneas que diga qué es, quién habla, cuánto dura, qué se editó y qué formato va a dónde.
\end{tcolorbox}

\begin{tcolorbox}[enhanced,colback=milcGris,colframe=milcGris,arc=2mm,boxrule=0pt,
  left=5mm,right=5mm,top=3.5mm,bottom=3.5mm,after skip=12pt]
{\sffamily\bfseries\fontsize{8}{10}\selectfont\textcolor{milcNegro!55}{ESTA GUÍA ES DE}}\\[3mm]
\begin{tabularx}{\linewidth}{X p{3.2cm} p{3.2cm}}
\linead{\linewidth} & \linead{3.0cm} & \linead{3.0cm}\\[-1mm]
{\fontsize{8.5}{10}\selectfont\textcolor{milcNegro!55}{Nombre completo}} &
{\fontsize{8.5}{10}\selectfont\textcolor{milcNegro!55}{Grupo}} &
{\fontsize{8.5}{10}\selectfont\textcolor{milcNegro!55}{Fecha}}\\
\end{tabularx}
\end{tcolorbox}

\vfill

\noindent\textcolor{milcLinea}{\rule{\linewidth}{1pt}}\\[2mm]
{\sffamily\bfseries\fontsize{9.5}{12}\selectfont Autor: Dr. Álvaro Cárdenas Orozco}\\[1mm]
{\sffamily\fontsize{8.5}{11}\selectfont\textcolor{milcNegro!55}{Metodología MILC · apertura ancestral · pensamiento computacional · triángulo Dussel–estoicismo–Floridi}}

\newpage

\section*{Apertura: Audio --- grabar, limpiar y exportar}

\begin{softbox}{milcGradeDark}{milcGradeSoft}{Saber ancestral}
En la sesión 3 conociste el saludo de Radio Pa'yumat. Hoy mira cómo llega esa voz a los resguardos. La emisora transmite en el 101.0 de la FM desde el cerro Munchique, en Santander de Quilichao, nueve horas al día. La sostienen cinco comunicadores indígenas coordinados por Abel Coicué (Aley, 2012). Los programas se graban y se emiten en nasa yuwe y en castellano: información, campañas, avisos de mingas, días de mercado, música. Antes de Pa'yumat hubo emisoras nasa en Toribío, en 1996, y en Jambaló, en 1997; Pa'yumat nació en 2000 como emisora de toda la zona. Lo que suena en el radio pasó por lo mismo que harás hoy: alguien grabó una voz, la limpió y la mandó por una antena. La \textbf{cara de exclusión}: en 2012 hubo atentados contra las antenas y torres; la señal perdió alcance y los equipos tuvieron que bajarse del cerro. Ningún editor de audio arregla una antena caída. Y la emisora depende del tiempo libre de niños, jóvenes y adultos que la sostienen. Hoy vas a grabar una voz, limpiarla y exportarla en dos formatos. En el norte del Cauca esa cadena se sostiene bajo amenaza.\par\smallskip{\footnotesize\textcolor{milcNegro!70}{\textbf{Fuente.} Aley, P. (2012, 22 de septiembre). Conozca la emisora de la resistencia indígena. \emph{El Tiempo}. · Tejido de Comunicación ACIN. (s.\,f.). \emph{Radio Pa'yumat}. Consultado el 7 de septiembre de 2026.}}
\end{softbox}

\begin{softbox}{milcTurquesa}{milcGris}{Saber contemporáneo}
Una pieza de audio se produce en tres fases. \textbf{Captura}: se graba con el micrófono del celular o del computador, en un lugar silencioso. El micrófono va a unos quince centímetros de la boca, y se hace una prueba de volumen antes. \textbf{Edición}: en Audacity, gratis en el computador, o en una app del celular, se hacen tres operaciones. Recortar: seleccionar el fragmento que sobra y eliminarlo. Ajustar el volumen: seleccionar todo y aplicar «Amplificar» hasta que se oiga parejo. Limpiar: quitar el silencio del inicio y del final, y reducir el ruido de fondo con «Reducción de ruido». \textbf{Exportación}: guardar en un formato que se pueda enviar. \textbf{MP3} comprime, pesa poco y sirve para WhatsApp y redes. \textbf{WAV} no comprime, pesa mucho más y sirve para guardar el original o seguir editando. Los errores típicos: volumen disparejo, ruido de ventilador, cortes bruscos sin fundido, y mandar un WAV pesado por donde iba un MP3.
\end{softbox}

\begin{softbox}{milcMostaza}{milcGris}{Pregunta-puente}
\textit{Pa'yumat graba en nasa yuwe, limpia la voz y la manda por una antena que puede caerse. Si grabas a tu abuela con el ventilador prendido, ¿qué parte de su voz se pierde y cuál de las tres operaciones lo arregla?}
\end{softbox}

\begin{softbox}{milcVerde}{milcGris}{Saber y saber hacer}
Al terminar podrás: (1) \textbf{identificar} qué sonó bien y qué sonó mal en una grabación tuya de treinta segundos; (2) \textbf{aplicar} recorte, volumen y limpieza en Audacity o una app; (3) \textbf{evaluar} qué formato va a cada destino; (4) \textbf{explicar} por qué el WAV se guarda y el MP3 se envía.
\end{softbox}



\small
\begin{tabularx}{\textwidth}{>{\bfseries}p{3.0cm}Y}
\rowcolor{milcGradePrimary!12}Autor & Dr. Álvaro Cárdenas Orozco\\
DBA & Producción multimedia ética y comportamiento responsable en entornos digitales (MEN, T\&I)\\
Referentes & Dussel (estética de la liberación) · Ley 1336 · Floridi (infoesfera)\\
Producto final & Una pieza de audio de uno a dos minutos, una entrevista corta a un familiar, una narración o un pódcast breve, grabada con el celular o el computador y editada con Audacity o una app del celular con tres operaciones: recorte, ajuste de volumen y limpieza de silencio o ruido. Exportada en dos formatos, MP3 y WAV, con una ficha de cinco líneas que diga qué es, quién habla, cuánto dura, qué se editó y qué formato va a dónde.\\
\end{tabularx}
\normalsize

\puente{En la sesión 4 editaste una imagen con permiso. Hoy editas una voz, y el permiso sigue: la voz de alguien también es suya. En la sesión 6 vas a entrar en un tema serio, el ciberacoso y cómo se denuncia.}

\milcestacion{milcGradeDark}{Ruta MILC: así vamos a aprender}
\begin{tabularx}{\textwidth}{>{\centering\arraybackslash}X>{\centering\arraybackslash}X>{\centering\arraybackslash}X>{\centering\arraybackslash}X}
\phasecard{1}{milcVerde}{milcGris}{Escucha\\[-1mm]\small Observo, escucho y pregunto.} &
\phasecard{2}{milcTurquesa}{milcGris}{Entender\\[-1mm]\small Sistematizo y ordeno.} &
\phasecard{3}{milcMagenta}{milcGris}{Hacer\\[-1mm]\small Construyo y aplico.} &
\phasecard{4}{milcVino}{milcGris}{Cierre\\[-1mm]\small Evalúo y reflexiono.}
\end{tabularx}


\puente{Antes de editar, vas a grabarte treinta segundos y escucharte. Tu propio oído te va a decir qué hay que limpiar.}

\milcestacion{milcVerde}{Estación 1 de 3 · Escucha}

\begin{softbox}{milcVerde}{milcGris}{Escena para escuchar}
\textbf{Actividad 1 · IDENTIFICA --- Treinta segundos de tu voz} (15 min · individual). (1) Con la grabadora del celular, graba treinta segundos hablando de cualquier cosa: cómo estuvo tu día, qué cambiarías del recreo. (2) Escúchate de inmediato con audífonos. (3) Anota qué sonó bien: claridad, energía, ritmo. (4) Anota qué sonó mal: ruido de fondo, volumen bajo, silencio largo al inicio, una palabra comida. (5) Para cada cosa que sonó mal, escribe si se arregla grabando de nuevo o editando.
\end{softbox}

{\sffamily\bfseries\fontsize{8}{10}\selectfont\textcolor{milcNegro!55}{MARCA CUANDO LO HAYAS HECHO}}
\milccheckrow{Tengo la grabación de treinta segundos.}
\milccheckrow{Anoté al menos una cosa que sonó bien y una que sonó mal.}
\milccheckrow{Para lo que sonó mal, dije si se arregla grabando o editando.}

\begin{infoband}{milcVerde}


\textbf{Lo que va al cuaderno (Actividad 1 · IDENTIFICA).} \textbf{Título:} «Actividad 1 --- Treinta segundos de tu voz». \textbf{Formato:} dos columnas, «sonó bien» y «sonó mal», y junto a cada problema la palabra «grabar» o «editar». \textbf{Extensión:} media página. \textbf{Sabes que terminaste cuando} sabes qué problema de tu grabación se arregla con cuál de las tres operaciones.
\end{infoband}

\puente{Ya sabes cómo suena tu voz grabada. Ahora aprendes las tres operaciones y, con tu pareja, limpias una grabación en Audacity o en la app.}

\milcestacion{milcTurquesa}{Estación 2 de 3 · Entender}

\begin{softbox}{milcTurquesa}{milcGris}{Lo que vas a entender}
\textbf{Actividad 2 · APLICA --- Las tres operaciones} (30 min · parejas). (1) Con tu pareja, abran en Audacity o en la app la grabación de treinta segundos de uno de los dos. (2) Recorten: seleccionen el silencio del inicio y del final y elimínenlo. (3) Ajusten el volumen: seleccionen todo y apliquen «Amplificar» hasta que la onda llene la pista sin tocar el borde. (4) Limpien: seleccionen dos segundos donde solo hay ruido de fondo, tomen el «perfil de ruido» y apliquen «Reducción de ruido» a todo. (5) Exporten como MP3 y como WAV, y comparen el peso de los dos archivos.
\end{softbox}

\milcpaso{1}{milcTurquesa}{\textbf{Recortar.} Clic y arrastrar sobre la onda para seleccionar; suprimir para eliminar. Se quita el silencio del inicio, el del final y las partes que sobran. Un fundido de un segundo al inicio y al final evita el corte brusco.}
\milcpaso{2}{milcTurquesa}{\textbf{Ajustar el volumen.} Seleccionar todo, menú Efecto, «Amplificar». La onda debe llenar la pista sin tocar el borde: si lo toca, se distorsiona. Si una parte quedó más baja, se selecciona solo esa y se amplifica aparte.}
\milcpaso{3}{milcTurquesa}{\textbf{Limpiar el ruido.} Se selecciona un pedazo donde solo hay ruido, ventilador o calle, y se toma el perfil. Luego «Reducción de ruido» a toda la pista. Poco reduce el ruido; mucho deja la voz metálica. Se prueba y se escucha.}
\milcpaso{4}{milcTurquesa}{\textbf{Cómo se hace, paso a paso.} (1) Abrir la grabación. (2) Recortar silencios. (3) Amplificar. (4) Perfil de ruido y reducción. (5) Fundidos. (6) Exportar MP3 y WAV. (7) Comparar el peso.}

\begin{drawbox}{milcTurquesa}{El audio, en una mirada}
\textbf{Captura:} lugar silencioso, micrófono a 15 cm, prueba de volumen. \\[1mm] \textbf{Recortar:} quitar lo que sobra. \\[1mm] \textbf{Amplificar:} volumen parejo sin tocar el borde. \\[1mm] \textbf{Reducir ruido:} perfil y aplicar. \\[1mm] \textbf{MP3:} liviano, para enviar. \textbf{WAV:} pesado, para guardar.
\end{drawbox}

\begin{softbox}{milcMostaza}{milcGris}{Errores comunes y su corrección}
(1) \textbf{Volumen disparejo}: una parte grita y otra no se oye. (2) \textbf{Ruido de fondo}: el ventilador o la calle encima de la voz. (3) \textbf{Corte brusco}: sin fundido, la pieza arranca o termina de golpe. (4) \textbf{Formato equivocado}: un WAV pesado por WhatsApp. (5) \textbf{Reducción de ruido exagerada}: la voz queda metálica.
\end{softbox}

\begin{infoband}{milcTurquesa}


\textbf{Lo que va al cuaderno (Actividad 2 · APLICA).} \textbf{Título:} «Actividad 2 --- Las tres operaciones». \textbf{Formato:} las tres operaciones con una línea sobre qué cambió en cada una, y el peso del MP3 y del WAV. \textbf{Extensión:} media página. \textbf{Sabes que terminaste cuando} la grabación limpia se oye pareja y sabes cuántas veces pesa más el WAV que el MP3.
\end{infoband}

\puente{Ya sabes limpiar. Toca grabar tu pieza de uno a dos minutos, editarla, exportarla en dos formatos y escribir la ficha.}

\milcestacion{milcMagenta}{Estación 3 de 3 · Hacer}

\begin{softbox}{milcMagenta}{milcGris}{Cómo vas a construir tu producto}
\textbf{Actividad 3 · EVALÚA --- Tu pieza en dos formatos} (25 min · individual). (1) Decide qué vas a grabar: una entrevista corta a un familiar, una narración o un pódcast breve, y pide permiso si graba a alguien. (2) Graba de uno a dos minutos en un lugar silencioso, con el micrófono a quince centímetros. (3) Aplica las tres operaciones y los fundidos. (4) Exporta en MP3 y en WAV y anota el peso de cada uno. (5) Escribe la ficha de cinco líneas. Pásale el MP3 a tu pareja para que lo oiga con audífonos y diga si notó algún corte o salto.
\end{softbox}

\milcpaso{1}{milcMagenta}{\textbf{Tu entrega tiene tres piezas.} (1) El MP3 y el WAV compartidos con enlace. (2) La ficha de cinco líneas: qué es, quién habla, cuánto dura, qué se editó, qué formato va a dónde. (3) Lo que dijo tu pareja al oír el MP3.}
\milcpaso{2}{milcMagenta}{\textbf{La voz de alguien es suya.} Si grabas a tu abuela o a un compañero, pides permiso antes y lo anotas, como con la imagen en la sesión 4. Y le dices para qué es y dónde va a sonar.}
\milcpaso{3}{milcMagenta}{\textbf{Qué formato va a dónde.} El WAV es el original: se guarda y no se manda. El MP3 es la copia liviana: se manda por WhatsApp, se sube a la red, se pone en la presentación. Mandar el WAV donde iba el MP3 es como mandar una caja donde cabía un sobre.}
\milcpaso{4}{milcMagenta}{\textbf{La prueba de los audífonos.} Tu pareja oye el MP3 sin saber dónde cortaste. Si no nota ningún salto, la edición está limpia. Si dice «ahí se oye un golpe», falta un fundido o un recorte quedó a mitad de palabra.}

\begin{drawbox}{milcMagenta}{Antes de entregar}
\checkbox Permiso pedido y anotado si grabé a alguien. \\[1mm] \checkbox Pieza de uno a dos minutos grabada en lugar silencioso. \\[1mm] \checkbox Recorte, amplificación y reducción de ruido aplicados. \\[1mm] \checkbox Fundido al inicio y al final. \\[1mm] \checkbox Exportado en MP3 y en WAV, con el peso de cada uno anotado. \\[1mm] \checkbox Ficha de cinco líneas. \\[1mm] \checkbox Mi pareja oyó el MP3 y no notó cortes ni saltos de volumen.
\end{drawbox}

\begin{softbox}{milcMagenta}{milcGris}{Plantilla de trabajo}
\textbf{Qué es.} \textit{Entrevista / narración / pódcast sobre …} \\[1mm] \textbf{Quién habla.} \textit{…, con permiso del …} \\[1mm] \textbf{Duración.} \textit{…:…} \\[1mm] \textbf{Qué se editó.} \textit{Recorté …, amplifiqué …, reduje el ruido de …} \\[1mm] \textbf{Formatos.} \textit{MP3 de … KB para …; WAV de … MB para guardar.}
\end{softbox}

\begin{infoband}{milcMagenta}


\textbf{Lo que va al cuaderno (Actividad 3 · EVALÚA).} \textbf{Título:} «Actividad 3 --- Tu pieza en dos formatos». \textbf{Formato:} los enlaces del MP3 y del WAV con su peso, la ficha de cinco líneas, la nota del permiso y lo que dijo tu pareja. \textbf{Extensión:} media página. \textbf{Sabes que terminaste cuando} tienes los dos archivos, la ficha llena y tu pareja no notó ningún corte.
\end{infoband}

\puente{Lo que entregas hoy: el MP3, el WAV y la ficha de cinco líneas. La prueba de calidad: tu pareja oye el MP3 con audífonos y no nota ningún corte ni cambio de volumen.}

\milcestacion{milcMostaza}{Producto de la sesión}
\textbf{Pieza de audio de uno a dos minutos con tres operaciones + MP3 y WAV + ficha de cinco líneas + prueba de los audífonos}.
\begin{softbox}{milcMostaza}{milcGris}{Criterios de calidad}
(1) La pieza dura entre uno y dos minutos y hay permiso anotado si habla alguien más. (2) Se aplicaron recorte, amplificación y reducción de ruido, con fundidos al inicio y al final. (3) Está exportada en MP3 y en WAV y el peso de cada uno está anotado. (4) La ficha dice qué es, quién habla, cuánto dura, qué se editó y qué formato va a dónde. (5) Tu pareja oyó el MP3 con audífonos y no notó cortes ni saltos de volumen. (6) Los dos archivos están compartidos con enlace.
\end{softbox}

\puente{Antes de cerrar, mira lo que hiciste desde cinco lados. Un minuto limpio y parejo vale más que cinco con ruido y saltos.}

\milcestacion{milcVino}{Evaluación liberadora · cinco dimensiones}

\begin{softbox}{milcVino}{milcGris}{Sentido de la evaluación}
Evaluar no es solo revisar una nota. Es reconocer qué aprendí, qué cambió en mí, qué emociones manejé, cómo me relaciono con la ciudad y la comunidad, y qué le dejaré a quien viene detrás.
\end{softbox}

\textbf{Mi reflexión liberadora en cinco dimensiones.} Completa cada frase iniciada:

\small
\begin{tabularx}{\textwidth}{>{\bfseries\RaggedRight\arraybackslash}p{3.4cm}Y>{\RaggedRight\arraybackslash}p{4.6cm}}
\rowcolor{milcVino}\color{white}Dimensión & \color{white}\bfseries Mi respuesta (frame starter) & \color{white}\bfseries Algo que descubrí\\
1. Personal & Lo que más cambió en mí fue \linead{6.5cm} & \linead{4.2cm}\\
2. Emocional & La emoción más fuerte fue \linead{2.5cm} y la regulé \linead{3cm} & \linead{4.2cm}\\
3. Ciudadana & En democracia esto importa porque \linead{6cm} & \linead{4.2cm}\\
4. Local (Cartago) & En mi barrio sería útil para \linead{6.5cm} & \linead{4.2cm}\\
5. Intergeneracional & A mi abuela/abuelo le diría \linead{6.5cm} & \linead{4.2cm}\\
\end{tabularx}
\normalsize

\begin{infoband}{milcVino}
\textbf{Para la próxima sustentación.} Vuelve a esta tabla en seis meses. Verás que las respuestas cambian — no porque lo aprendido cambie, sino porque tú cambias. La evaluación liberadora es una conversación contigo a través del tiempo.
\end{infoband}


\puente{Para terminar, tres ideas sobre grabar la voz de otro, contadas con palabras de esta guía: Dussel sobre el silencio del que escucha, Epicteto sobre usar las menos palabras, y Floridi sobre las tecnologías que ahorran tiempo y luego lo matan.}

\milcestacion{milcGradeDark}{Tres ideas para cerrar · Dussel, un estoico y Floridi}

\begin{softbox}{milcVino}{milcGris}{Enrique Dussel · lente del nosotros}
\textit{El respeto es el silencio de quien todo tiene que escuchar, porque todavía no sabe nada del otro.}

{\footnotesize\itshape\color{milcNegro!70}--- idea tomada de Enrique Dussel, contada con palabras de esta guía. No es una cita textual.}

Grabar a tu abuela es callarte quince centímetros del micrófono y dejar que hable. Y editar su voz es quitarle el ruido, no las palabras que a ti te parecieron de más.

\textbf{¿Qué le recorté a la voz de alguien porque yo lo habría dicho distinto?}
\end{softbox}
\linead{\linewidth}\\[1mm]
\linead{\linewidth}

\begin{softbox}{milcMostaza}{milcGris}{Estoicismo · Epicteto · cuidado interior}
\textit{Guarda silencio cuanto puedas; di solo lo necesario y con las menos palabras posibles.}

{\footnotesize\itshape\color{milcNegro!70}--- idea tomada de Epicteto, contada con palabras de esta guía. No es una cita textual.}

Una pieza de un minuto limpia dice más que cinco con silencios y ruido. Recortar es la forma de que lo necesario se oiga. Pa'yumat tiene nueve horas al día y aun así cada aviso es corto.

\textbf{¿Cuántos segundos de mi pieza sobraban y me costó cortar?}
\end{softbox}
\linead{\linewidth}\\[1mm]
\linead{\linewidth}

\begin{softbox}{milcTurquesa}{milcGris}{Luciano Floridi · lente de la infoesfera}
\textit{Las tecnologías se usan primero para ahorrar tiempo y después para matarlo.}

{\footnotesize\itshape\color{milcNegro!70}--- idea tomada de Luciano Floridi, contada con palabras de esta guía. No es una cita textual.}

Un MP3 de un minuto que alguien oye en el bus le ahorra tiempo. Un WAV de veinte megas que no descarga, o un audio de cinco minutos con ruido, se lo mata. Elegir formato y duración es decidir cuál de las dos cosas haces.

\textbf{¿Qué audio mandé esta semana que le hizo perder tiempo a quien lo recibió?}
\end{softbox}
\linead{\linewidth}\\[1mm]
\linead{\linewidth}

\begin{infoband}{milcNegro}
\textbf{Nota para el docente.} Las tres ideas están contadas con palabras de esta guía a partir del pensamiento de cada autor; no son citas textuales. De 9.º en adelante se trabajan con la obra y su referencia.
\end{infoband}

\milcestacion{milcVino}{Mi autoevaluación · marca una casilla por fila}

{\small
\setlength{\extrarowheight}{2.2pt}
\begin{tabularx}{\textwidth}{>{\RaggedRight\arraybackslash\bfseries}p{3.0cm}YYY}
\rowcolor{milcVino}\color{white}\bfseries Criterio & \color{white}\bfseries Lo logré & \color{white}\bfseries En proceso & \color{white}\bfseries Necesito apoyo\\[.6mm]
Comprensión del tema & \milcopcion{Explico las ideas con mis palabras.} & \milcopcion{Reconozco las ideas, falta precisión.} & \milcopcion{Necesito repasar lo básico.}\\[.8mm]
\rowcolor{milcGris}Pensamiento computacional & \milcopcion{Usé los pilares al armar mi producto.} & \milcopcion{Usé algunos pilares.} & \milcopcion{Necesito releer la estación 2.}\\[.8mm]
Aplicación y producto & \milcopcion{Mi producto cumple los criterios.} & \milcopcion{Está armado, con ajustes pendientes.} & \milcopcion{Aún está en construcción.}\\[.8mm]
\rowcolor{milcGris}Reflexión 5 dimensiones & \milcopcion{Respondí cada una con honestidad.} & \milcopcion{Respondí algunas con detalle.} & \milcopcion{Mis respuestas son muy generales.}\\[.8mm]
Triángulo de pensamiento & \milcopcion{Conecté las 3 voces con mi trabajo.} & \milcopcion{Conecté al menos una.} & \milcopcion{No las relaciono con mi proceso.}\\[.8mm]
\end{tabularx}}

\vspace{3mm}
\textbf{Hoy entendí que:}\\[1mm]
\linead{\linewidth}\\[1mm]
\linead{\linewidth}

\textbf{Mi compromiso para la próxima sesión es:}\\[1mm]
\linead{\linewidth}\\[1mm]
\linead{\linewidth}

\begin{softbox}{milcMostaza}{milcGris}{Cita de despedida · Marco Aurelio}
\textit{``El obstáculo en el camino se vuelve el camino. Lo que estorba al camino se vuelve camino.''}\\[1mm]
Cada error de hoy, bien observado, se convierte en pieza del camino que pisarás mañana.
\end{softbox}

\small
\begin{tabularx}{\textwidth}{>{\bfseries}p{3.6cm}Y}
\rowcolor{milcGradeDark!12}Nombre completo: & \linead{10cm}\\
Fecha: & \linead{4cm} \quad Curso/grupo: \linead{4cm}\\
Mi firma: & \linead{9cm}\\
\end{tabularx}
\normalsize

\begin{infoband}{milcGradeDark}
\textbf{Para la próxima sesión.} Llega con tu MP3, tu WAV y la ficha. En la sesión 6 vas a tratar un tema serio: el ciberacoso, cómo se reconoce y cómo se denuncia, con la Línea 141 del ICBF. Lo que aprendiste sobre el permiso para usar la imagen y la voz de alguien es la base de esa conversación.
\end{infoband}

\end{document}
